Multilingual retrieval benchmarks are now common in general-domain evaluation, but legal retrieval remains comparatively underrepresented, especially in settings where relevance must hold across several language realizations of the same document. Existing multilingual retrieval benchmarks often focus on web, encyclopedic, or question-answering data whose relevance structure is simpler than that of legal corpora. By contrast, legal retrieval requires stable document identity, careful handling of long structured texts, and supervision that reflects substantive legal information needs rather than surface lexical overlap.

Within legal NLP, prior work has emphasized tasks such as statute or case retrieval, legal entailment, citation prediction, and summarization. These tasks have shown that legal text is unusually sensitive to document structure, formulaic phrasing, and domain-specific terminology. They also highlight a recurring tension between faithful extraction and realistic user-facing information needs. Our benchmark is positioned in that gap: it is retrieval-oriented, multilingual, and document-level, but it still preserves an auditable trace from source text to generated supervision.

The construction strategy is also related to recent work that uses large language models to build retrieval datasets through synthetic question--answer generation and validation. That literature shows that synthetic supervision can be useful when human annotation is expensive, but it also makes clear that generation quality depends heavily on prompt design, filtering, and domain-specific validation. In legal text, those concerns are amplified by provision-led wording, inventory-style questions, and clause-heavy prompts that look superficially valid while remaining weak retrieval queries. The paper therefore situates itself at the intersection of multilingual retrieval benchmarking, legal information retrieval, and synthetic supervision for retrieval evaluation.

Formal citation support and a more complete literature comparison will be added once the paper bibliography is assembled.
